
\lstset{
  basicstyle=\ttfamily\small,
  breaklines=true,
  columns=fullflexible
}

\section{{\tt elegant} Command-file Mode for Emacs}

\subsection{Introduction}

This document describes the installation and features of a custom
\texttt{elegant-mode} for GNU Emacs.
The mode provides structured editing support for
\texttt{.ele} input files used by {\tt elegant}.

The mode includes:
\begin{itemize}[noitemsep]
\item Syntax highlighting for namelists and qualifiers
\item Auto-completion of namelist names
\item Auto-completion of qualifiers within namelist blocks
\item Automatic insertion of \texttt{\&end}
\item Automatic indentation (4 spaces for qualifier lines)
\item Automatic activation for \texttt{.ele} files
\end{itemize}

No external Emacs packages are required.

\subsection{Features}

\subsubsection{Namelist Completion}

For illustration, we use the \verb|run_setup| command. The same techniques apply to all
commands.
When typing:
\begin{lstlisting}
&run_
\end{lstlisting}
then pressing \texttt{TAB} completes from the list of known {\tt elegant}
namelists (e.g. \texttt{run_setup}).

Typing
\begin{lstlisting}
&run_setup
\end{lstlisting}
then pressing \texttt{RET} produces:
\begin{lstlisting}
&run_setup
    <cursor here>
&end
\end{lstlisting}
The cursor is positioned inside the block, ready for qualifier entry.

\subsubsection{Qualifier Completion}

Inside a namelist block, pressing \texttt{TAB} completes qualifier names
specific to that namelist:
\begin{lstlisting}
&run_setup
    latt<TAB>
&end
\end{lstlisting}
Completes to:
\begin{lstlisting}
    lattice =
\end{lstlisting}

\subsection{Indentation Rules}

\begin{itemize}[noitemsep]
\item Namelist headers begin at column 0.
\item \texttt{\&end} begins at column 0.
\item Qualifier lines are indented 4 spaces.
\item Pressing \texttt{RET} maintains indentation.
\end{itemize}

\subsection{Installation}

\subsubsection{Directory Structure}

Create the following directory structure in your home directory:
\begin{lstlisting}
/home/<username>/.emacs
/home/<username>/.emacs.d/elegant/
\end{lstlisting}

The files required are:
\begin{lstlisting}
~/.emacs
~/.emacs.d/elegant/elegant-mode.el
~/.emacs.d/elegant/elegant-namelists.el
\end{lstlisting}

\subsubsection{Step 1: Generate Namelist Tables}

This part is usually done by the {\tt elegant} developers. If you install {\tt elegant} from source,
you may also wish to do this step. Use the provided Python script (in {\tt elegantTools}):
\begin{lstlisting}
cat *.nl > elegant-namelists.txt
python3 gen_elegant_elisp.py elegant-namelists.txt \
  > ~/.emacs.d/elegant/elegant-namelists.el
\end{lstlisting}

This creates:
\begin{itemize}[noitemsep]
\item \texttt{elegant-namelist-commands}
\item \texttt{elegant-namelist-qualifiers}
\end{itemize}

These are required for completion.

\subsubsection{Step 2: Install elegant-mode.el}

Copy the provided \texttt{elegant-mode.el} file into:
\begin{lstlisting}
~/.emacs.d/elegant/
\end{lstlisting}

\subsubsection{Step 3: Configure Emacs}

Edit or create:
\begin{lstlisting}
/home/<username>/.emacs
\end{lstlisting}

\noindent
Add:\\
{\tt
(add-to-list \textquotesingle load-path "~/.emacs.d/elegant") \\
(load "elegant-namelists.el" nil t) \\
(load "elegant-mode.el" nil t)\\
(add-to-list \textquotesingle auto-mode-alist \textquotesingle ("\textbackslash\textbackslash .ele\textbackslash\textbackslash\textquotesingle " . elegant-mode)) \\
(setq completion-cycle-threshold 1) \\
(setq tab-always-indent \textquotesingle complete) \\
(define-key global-map (kbd "TAB") \#\textquotesingle completion-at-point) \\
}

\subsubsection{Step 4: Restart Emacs}

Start Emacs normally:

\begin{lstlisting}
emacs run.ele
\end{lstlisting}

The mode line should display {\tt {\tt elegant}}.
To confirm correct installation:

\begin{enumerate}[noitemsep]
\item Open a \texttt{.ele} file.
\item Type \texttt{\&run_setup} and press \texttt{RET}.
\item Confirm that \texttt{\&end} is inserted automatically.
\item Type part of a qualifier and press \texttt{TAB}.
\item Confirm completion suggestions appear.
\end{enumerate}

\subsection{Customization}

Indentation width may be adjusted:
\begin{lstlisting}
(setq elegant-indent-offset 4)
\end{lstlisting}

This variable can be changed to any desired indentation width.

\subsection{Troubleshooting}

\subsubsection{Mode Not Activating}

Verify:
\begin{lstlisting}
M-x describe-variable RET auto-mode-alist RET
\end{lstlisting}

Ensure:
\begin{lstlisting}
("\\.ele\\' " . elegant-mode)
\end{lstlisting}
is present.

\subsubsection{Completion Not Working}

Verify:
\begin{lstlisting}
M-: elegant-namelist-commands RET
\end{lstlisting}
returns a list of namelists.

If it returns \texttt{nil}, confirm that
\texttt{elegant-namelists.el} was generated and loaded.

\subsection{Conclusion}

\texttt{elegant-mode} provides a lightweight, self-contained editing
environment for {\tt elegant} input files with:

\begin{itemize}[noitemsep]
\item Structured completion
\item Automatic block management
\item Clean indentation
\item Zero external dependencies
\end{itemize}

It integrates seamlessly into standard Emacs installations and requires only three files.


